\paragraph{Overall verdict.}
\emph{Partially reproduced}. With today's \asan{}, the average run-time claims (C1,
C2) and the stack and binary-size claims hold for programs comparable to the
paper's, while the 3.37$\times$ memory increase does not: we measure
\OurMemTotal$\times$ in peak resident memory, and the paper's VmPeak metric no
longer applies.

\paragraph{Lessons learned.}
\begin{itemize}
\item \emph{Settings matter as much as the mechanism.} Run-time options
(quarantine, allocation stack traces, redzone size) changed our slowdown by
up to \CtrlDefDelta{} on average and our memory increase by a factor of
almost three. A reproduction has to copy these settings, not just the
compiler flag.
\item \emph{Metrics age.} VmPeak was a sensible memory metric in 2012; with
today's 20\,TB shadow reservation it measures nothing. We had to justify a
replacement (VmHWM) and record both.
\item \emph{Workload composition drives the averages.} The slowdown is
largely explained by allocation rate. Our 11-program average is
\GapRWAllRel{} above the paper's headline number, but the like-for-like
comparison with the same seven SPEC programs agrees to within
\GapRWMapRel. A headline average is only meaningful together with the
program set it was computed on.
\item \emph{Cloud VMs are noisy.} Steal time of up to \StealMaxPct\% of a run made
interleaved repetitions, medians and steal accounting necessary before any
average could be trusted.
\end{itemize}
